\documentclass[12pt,a4paper]{article}
\usepackage[utf8]{inputenc}
\usepackage[spanish]{babel}
\usepackage{amsmath}
\usepackage{amsfonts}
\usepackage{amssymb}
\usepackage{graphicx}
\usepackage[left=2.5cm,right=2.5cm,top=3cm,bottom=3cm]{geometry}
\usepackage{xcolor}
\usepackage{listings}
\usepackage{tcolorbox}
\usepackage{hyperref}

% Configuración de colores para código
\definecolor{codegreen}{rgb}{0,0.6,0}
\definecolor{codegray}{rgb}{0.5,0.5,0.5}
\definecolor{codepurple}{rgb}{0.58,0,0.82}
\definecolor{backcolour}{rgb}{0.95,0.95,0.92}

\lstdefinestyle{mystyle}{
    backgroundcolor=\color{backcolour},   
    commentstyle=\color{codegreen},
    keywordstyle=\color{magenta},
    numberstyle=\tiny\color{codegray},
    stringstyle=\color{codepurple},
    basicstyle=\ttfamily\footnotesize,
    breakatwhitespace=false,         
    breaklines=true,                 
    captionpos=b,                    
    keepspaces=true,                 
    numbers=left,                    
    numbersep=5pt,                  
    showspaces=false,                
    showstringspaces=false,
    showtabs=false,                  
    tabsize=2
}
\lstset{style=mystyle}

\title{\textbf{Guía de Laboratorio 3}\\ \Large Laboratorio Virtual: Doppler Nexus}
\author{Andrés Jaramillo - Andrea Zapata - Víctor Gonzales - Sergio Carmona - Daniela Nieto}
\date{\today}

\begin{document}

\maketitle

\section{Introducción}
Este proyecto consiste en una simulación interactiva avanzada diseñada para explorar los principios fundamentales del Efecto Doppler. A través de un motor gráfico en 2D, se visualiza cómo la velocidad de una fuente sonora comprime o expande los frentes de onda, alterando la frecuencia percibida por un observador en reposo o en movimiento.

\section{Fundamento Teórico}
El efecto Doppler es el cambio aparente en la frecuencia de una onda causado por el movimiento relativo entre la fuente de la onda y su observador.

\begin{tcolorbox}[colback=purple!5!white,colframe=purple!75!black,title=Ecuaciones del Efecto Doppler]
\textbf{General:} 
\[ f' = f \left( \frac{v \pm v_o}{v \mp v_s} \right) \]
Donde:
\begin{itemize}
    \item $v$: Velocidad del sonido en el medio.
    \item $v_o$: Velocidad del observador.
    \item $v_s$: Velocidad de la fuente.
\end{itemize}
\end{tcolorbox}

\section{Características de la Simulación}
\begin{itemize}
    \item \textbf{Propagación Dinámica:} Visualización de crestas de onda como círculos concéntricos.
    \item \textbf{Número de Mach:} Simulación de regímenes subsónicos, sónicos ($M=1$) y supersónicos ($M>1$).
    \item \textbf{Cono de Mach:} Formación de la onda de choque y explosión sónica (sonic boom).
\end{itemize}

\section{Algoritmo de Renderizado (Extracto)}
\begin{lstlisting}[language=JavaScript]
// Manejo de ondas en el canvas
function renderLoop() {
    ctx.clearRect(0, 0, width, height);
    
    // Nueva onda cada intervalo
    if (ticks % interval === 0) {
        ondas.push({ x: source.x, y: source.y, r: 0 });
    }
    
    // Dibujo y actualizacion de radios
    ondas.forEach((o, i) => {
        o.r += speedSound;
        ctx.beginPath();
        ctx.arc(o.x, o.y, o.r, 0, 2 * Math.PI);
        ctx.stroke();
        
        // Limpieza de memoria
        if (o.r > maxRadius) ondas.splice(i, 1);
    });
    
    requestAnimationFrame(renderLoop);
}
\end{lstlisting}

\section{Análisis de Resultados}
Al utilizar el simulador, se observa que:
\begin{enumerate}
    \item Cuando la fuente se acerca al observador, los frentes de onda se agrupan ($f' > f$).
    \item Cuando la fuente se aleja, los frentes se distancian ($f' < f$).
\end{enumerate}

\begin{figure}[h]
    \centering
    % \includegraphics[width=0.8\textwidth]{screenshot_doppler.png}
    \caption{Visualización del efecto Doppler en el simulador Nexus (Espacio para captura).}
\end{figure}

\end{document}
